\section{Introduction}

Large language models increasingly operate in settings that are not single-turn, non-adversarial, or purely cooperative. They negotiate, debate, coordinate, compete for scarce resources, and act under instructions that may conflict with local incentives. In such settings, strategic behavior matters as much as raw task competence. A model that reasons well but systematically misrepresents information under pressure poses a different alignment concern than a model that simply makes mistakes.

Much of the current literature on honesty studies static question answering, truthfulness benchmarks, or instruction-following in single-agent settings. Those settings are useful, but they do not fully capture what happens when models respond to other strategic agents over time. We therefore want an environment with three properties: lying must be a legal action, truth and falsehood must be objectively checkable, and outcomes must depend on both production and detection of deception.

Bullshit provides those properties with relatively low task overhead. Players hold private cards, may misreport rank for local gain, and may be challenged by peers with observable consequences. For a TMLR audience, the value of the benchmark is not that it is a card game. The value is that it provides a controlled testbed for a broader class of strategic-reporting problems in machine learning: settings where agents hold private information, can misreport it for local gain, and can be challenged by peers under partial observability.

This paper asks three concrete questions. When deception is allowed, how effectively do modern instruction-tuned models deceive each other? When they believe opponents are constrained to play honestly, do they show any restraint? When explicitly instructed not to lie, do they continue to do so anyway? We answer these questions through Bullshit-Bench, a reproducible benchmark release and pilot dataset for strategic misrepresentation, challenge behavior, and instruction compliance in multi-agent LLM play.

\paragraph{Contributions.}
We make four contributions. First, we introduce a benchmark release for competitive deception evaluation in a game where lies, challenges, and outcomes are all directly measurable. Second, we provide a reproducible protocol based on seeded execution, a public environment API, frozen prompt versions, cohort filtering, and uncapped winner-terminated runs. Third, we ship an end-to-end artifact pipeline from gameplay logs to CSV exports, figures, release manifests, and replayable qualitative examples. Fourth, we report an empirical pilot study in which each headline result is tied to explicit evidence files rather than to anecdotal gameplay.
